% Options for packages loaded elsewhere
% Options for packages loaded elsewhere
\PassOptionsToPackage{unicode}{hyperref}
\PassOptionsToPackage{hyphens}{url}
\PassOptionsToPackage{dvipsnames,svgnames,x11names}{xcolor}
%
\documentclass[
  11pt,
]{article}
\usepackage{xcolor}
\usepackage[margin=1in]{geometry}
\usepackage{amsmath,amssymb}
\setcounter{secnumdepth}{-\maxdimen} % remove section numbering
\usepackage{iftex}
\ifPDFTeX
  \usepackage[T1]{fontenc}
  \usepackage[utf8]{inputenc}
  \usepackage{textcomp} % provide euro and other symbols
\else % if luatex or xetex
  \usepackage{unicode-math} % this also loads fontspec
  \defaultfontfeatures{Scale=MatchLowercase}
  \defaultfontfeatures[\rmfamily]{Ligatures=TeX,Scale=1}
\fi
\usepackage{lmodern}
\ifPDFTeX\else
  % xetex/luatex font selection
\fi
% Use upquote if available, for straight quotes in verbatim environments
\IfFileExists{upquote.sty}{\usepackage{upquote}}{}
\IfFileExists{microtype.sty}{% use microtype if available
  \usepackage[]{microtype}
  \UseMicrotypeSet[protrusion]{basicmath} % disable protrusion for tt fonts
}{}
\usepackage{setspace}
\makeatletter
\@ifundefined{KOMAClassName}{% if non-KOMA class
  \IfFileExists{parskip.sty}{%
    \usepackage{parskip}
  }{% else
    \setlength{\parindent}{0pt}
    \setlength{\parskip}{6pt plus 2pt minus 1pt}}
}{% if KOMA class
  \KOMAoptions{parskip=half}}
\makeatother
% Make \paragraph and \subparagraph free-standing
\makeatletter
\ifx\paragraph\undefined\else
  \let\oldparagraph\paragraph
  \renewcommand{\paragraph}{
    \@ifstar
      \xxxParagraphStar
      \xxxParagraphNoStar
  }
  \newcommand{\xxxParagraphStar}[1]{\oldparagraph*{#1}\mbox{}}
  \newcommand{\xxxParagraphNoStar}[1]{\oldparagraph{#1}\mbox{}}
\fi
\ifx\subparagraph\undefined\else
  \let\oldsubparagraph\subparagraph
  \renewcommand{\subparagraph}{
    \@ifstar
      \xxxSubParagraphStar
      \xxxSubParagraphNoStar
  }
  \newcommand{\xxxSubParagraphStar}[1]{\oldsubparagraph*{#1}\mbox{}}
  \newcommand{\xxxSubParagraphNoStar}[1]{\oldsubparagraph{#1}\mbox{}}
\fi
\makeatother


\usepackage{longtable,booktabs,array}
\usepackage{calc} % for calculating minipage widths
% Correct order of tables after \paragraph or \subparagraph
\usepackage{etoolbox}
\makeatletter
\patchcmd\longtable{\par}{\if@noskipsec\mbox{}\fi\par}{}{}
\makeatother
% Allow footnotes in longtable head/foot
\IfFileExists{footnotehyper.sty}{\usepackage{footnotehyper}}{\usepackage{footnote}}
\makesavenoteenv{longtable}
\usepackage{graphicx}
\makeatletter
\newsavebox\pandoc@box
\newcommand*\pandocbounded[1]{% scales image to fit in text height/width
  \sbox\pandoc@box{#1}%
  \Gscale@div\@tempa{\textheight}{\dimexpr\ht\pandoc@box+\dp\pandoc@box\relax}%
  \Gscale@div\@tempb{\linewidth}{\wd\pandoc@box}%
  \ifdim\@tempb\p@<\@tempa\p@\let\@tempa\@tempb\fi% select the smaller of both
  \ifdim\@tempa\p@<\p@\scalebox{\@tempa}{\usebox\pandoc@box}%
  \else\usebox{\pandoc@box}%
  \fi%
}
% Set default figure placement to htbp
\def\fps@figure{htbp}
\makeatother





\setlength{\emergencystretch}{3em} % prevent overfull lines

\providecommand{\tightlist}{%
  \setlength{\itemsep}{0pt}\setlength{\parskip}{0pt}}



 


\makeatletter
\@ifpackageloaded{caption}{}{\usepackage{caption}}
\AtBeginDocument{%
\ifdefined\contentsname
  \renewcommand*\contentsname{Table of contents}
\else
  \newcommand\contentsname{Table of contents}
\fi
\ifdefined\listfigurename
  \renewcommand*\listfigurename{List of Figures}
\else
  \newcommand\listfigurename{List of Figures}
\fi
\ifdefined\listtablename
  \renewcommand*\listtablename{List of Tables}
\else
  \newcommand\listtablename{List of Tables}
\fi
\ifdefined\figurename
  \renewcommand*\figurename{Figure}
\else
  \newcommand\figurename{Figure}
\fi
\ifdefined\tablename
  \renewcommand*\tablename{Table}
\else
  \newcommand\tablename{Table}
\fi
}
\@ifpackageloaded{float}{}{\usepackage{float}}
\floatstyle{ruled}
\@ifundefined{c@chapter}{\newfloat{codelisting}{h}{lop}}{\newfloat{codelisting}{h}{lop}[chapter]}
\floatname{codelisting}{Listing}
\newcommand*\listoflistings{\listof{codelisting}{List of Listings}}
\makeatother
\makeatletter
\makeatother
\makeatletter
\@ifpackageloaded{caption}{}{\usepackage{caption}}
\@ifpackageloaded{subcaption}{}{\usepackage{subcaption}}
\makeatother
\usepackage{bookmark}
\IfFileExists{xurl.sty}{\usepackage{xurl}}{} % add URL line breaks if available
\urlstyle{same}
\hypersetup{
  pdftitle={Latent Pathways and Post-Disruption Reorganization in Mexican Cartel Networks},
  colorlinks=true,
  linkcolor={blue},
  filecolor={Maroon},
  citecolor={Blue},
  urlcolor={Blue},
  pdfcreator={LaTeX via pandoc}}


\title{Latent Pathways and Post-Disruption Reorganization in Mexican
Cartel Networks}
\usepackage{etoolbox}
\makeatletter
\providecommand{\subtitle}[1]{% add subtitle to \maketitle
  \apptocmd{\@title}{\par {\large #1 \par}}{}{}
}
\makeatother
\subtitle{Project Proposal --- Networks I: Concepts and Algorithms}
\author{}
\date{}
\begin{document}
\maketitle


\setstretch{1.15}
\subsection{What I Want to Do}\label{what-i-want-to-do}

Kingpin strategies assume that removing central cartel actors fragments
criminal organizations. Empirical evidence says otherwise: removals in
Mexico triggered violence spikes and reconfiguration, not collapse
(Calderón et al., 2015; Duijn et al., 2014). I argue that cartel
networks contain \emph{latent relational pathways} --- shared enemies,
brokerage alternatives, overlapping neighborhoods --- that become
activated when central actors are removed. The question is not only
whether disruption fragments a network, but \textbf{whether pre-existing
latent structure conditions how it reorganizes}.

Using the \textbf{BACRIM2020} dataset (Prieto-Curiel \& Campedelli,
2023) --- 150 criminal factions in Mexico with weighted alliance and
rivalry ties --- I model cartel relations as a \textbf{signed network}
and test:

\begin{quote}
\textbf{H1:} Pre-existing latent relational pathways condition the
likely form of post-disruption reorganization after the removal of
structurally central actors.
\end{quote}

This is policy-relevant: if ``enemy-of-my-enemy'' dynamics generate new
coalitions after kingpin removals, understanding the latent structure
\emph{before} intervention enables smarter targeting. Two real-world
anchors motivate the design: (1) the post-2017 Sinaloa fragmentation
after El Chapo's extradition, and (2) the recent death of ``El Mencho''
(CJNG), whose node (N10025) is the single most connected in the dataset
--- 28 alliances and 29 rivalries --- making it the strongest test case
for whether latent pathways predict reconfiguration.

\subsection{How I Plan to Do It}\label{how-i-plan-to-do-it}

The analysis has six steps:

\textbf{1. Signed Network Construction.} Build alliance (+1) and rivalry
(−1) layers on the same 150 nodes. Analyze them separately and jointly.

\textbf{2. Central Actor Identification.} Compute degree, betweenness,
and eigenvector centrality plus structural-hole measures. Focus on
brokers --- actors whose removal exposes dormant coordination
alternatives (Morselli, 2009; Bright et al., 2015).

\textbf{3. Latent Pathway Inference (Pre-Disruption).} For each
non-adjacent dyad, compute a composite latent score from: Jaccard
similarity in alliance/rivalry neighborhoods, Adamic-Adar overlap
(Adamic \& Adar, 2003; Liben-Nowell \& Kleinberg, 2007), community
co-membership, and short indirect path length. Shared rivals receive
extra weight --- in criminal contexts, mutual enemies create incentives
for coordination even without a visible alliance (Everton, 2012).

\textbf{4. Node Removal Simulation.} Remove central nodes (e.g., CJNG,
Sinaloa) and measure fragmentation, component size, bridge loss, and
signed triad balance (Albert et al., 2000; Wandelt et al., 2018). This
approximates targeted intervention effects.

\textbf{5. Counterfactual Reconstruction.} Add new ties \emph{only}
where pre-removal latent scores are high (Leskovec et al., 2010). Assess
whether plausible recombination reduces fragmentation, restores
brokerage, or generates new coalitional blocs.

\textbf{6. Validation via Link Prediction.} Withhold a random subset of
observed ties, recalculate latent scores from the remaining structure,
and check whether withheld ties rank higher than unrelated dyads.
Evaluate with AUC and precision-at-k.

The project is implemented in \textbf{Python} (\texttt{networkx},
\texttt{numpy}, \texttt{pandas}, \texttt{matplotlib}) and delivered as a
\textbf{Quarto} report.

\subsection{Challenges and How to Overcome
Them}\label{challenges-and-how-to-overcome-them}

\begin{longtable}[]{@{}
  >{\raggedright\arraybackslash}p{(\linewidth - 2\tabcolsep) * \real{0.5000}}
  >{\raggedright\arraybackslash}p{(\linewidth - 2\tabcolsep) * \real{0.5000}}@{}}
\toprule\noalign{}
\begin{minipage}[b]{\linewidth}\raggedright
Challenge
\end{minipage} & \begin{minipage}[b]{\linewidth}\raggedright
Mitigation
\end{minipage} \\
\midrule\noalign{}
\endhead
\bottomrule\noalign{}
\endlastfoot
\textbf{Cross-sectional data} --- cannot observe actual post-disruption
dynamics. & Frame results as \emph{counterfactual plausibility}, not
causal claims. Supplement with qualitative comparison to documented
Sinaloa/CJNG reconfiguration. \\
\textbf{Partial observability} --- latent edges are unobserved by
definition. & The link-prediction validation (Step 6) demonstrates the
score recovers artificially hidden edges, establishing it as a credible
proxy. \\
\textbf{Threshold sensitivity} --- results may depend on how many ties
are added. & Sensitivity analysis across multiple thresholds; report how
fragmentation recovery and coalition formation vary. \\
\textbf{Signed network complexity} --- standard metrics assume unsigned
graphs. & Compute alliance- and rivalry-specific versions of each metric
separately; use signed triad balance (Cartwright \& Harary, 1956) to
validate structural coherence. \\
\textbf{Small N (\textasciitilde150 nodes)} --- limited statistical
power. & Use rank-based metrics (AUC) robust at smaller scales and
repeated random splits. \\
\end{longtable}

\subsection{Key References}\label{key-references}

Adamic \& Adar (2003). \emph{Social Networks} 25(3). --- Albert, Jeong
\& Barabási (2000). \emph{Nature} 406. --- Bright et al.~(2015).
\emph{J. Contemp. Crim. Justice} 31(3). --- Calderón et al.~(2015).
\emph{J. Conflict Resolution} 59(8). --- Cartwright \& Harary (1956).
\emph{Psych. Review} 63(5). --- Duijn, Kashirin \& Sloot (2014).
\emph{Sci. Reports} 4. --- Everton (2012). \emph{Disrupting Dark
Networks}, CUP. --- Leskovec, Huttenlocher \& Kleinberg (2010).
\emph{WWW 2010}. --- Liben-Nowell \& Kleinberg (2007). \emph{JASIST}
58(7). --- Lü \& Zhou (2011). \emph{Physica A} 390(6). --- Morselli
(2009). \emph{Inside Criminal Networks}, Springer. --- Prieto-Curiel \&
Campedelli (2023). BACRIM2020 dataset. --- Wandelt et al.~(2018).
\emph{Sci. Reports} 8.




\end{document}
